% auto-generated from vault memos --- do not edit by hand
\begin{tcolorbox}[colback=gray!4,colframe=black!35,boxrule=0.4pt,left=3pt,right=3pt,top=2pt,bottom=2pt,arc=1pt,title={\small\textbf{S3M1} --- earlier BPMN contact unlocks MSP start}]
\footnotesize \textbf{Case:} S3. \textbf{Codes:} S3-C04, S3-C05, S3-C06, S3-C17. \textbf{Turns:} T013, T017, T019, T043.

\textit{Idea.} S3 encountered BPMN in two brief earlier courses before MSP. That prior contact let him skip basic assembly and start exercises as soon as they appeared, while a peer without that history needed more kick-start help. Crash-course conditions (about one week of BPMN, then the course moved on) built familiarity but not depth; MSP then added specific symbols and terminology on top.

\textit{Why it matters.} For STGT, repeated socio-technical exposure---notation plus tool plus short practical windows---changed what ``starting to model'' felt like in the main modelling course. Prior contact with the same tool/notation pair (Bizagi + BPMN) reduced dependence on waiting for a board example.
\end{tcolorbox}
\vspace{2pt}
\begin{tcolorbox}[colback=gray!4,colframe=black!35,boxrule=0.4pt,left=3pt,right=3pt,top=2pt,bottom=2pt,arc=1pt,title={\small\textbf{S8M1} --- modelling arrived as a project requirement first}]
\footnotesize \textbf{Case:} S8. \textbf{Codes:} S8-C01, S8-C03, S8-C04. \textbf{Turns:} T003, T011, T013.

\textit{Idea.} At ISEL, S8 first had to produce UML/RE diagrams in the bachelor's project without knowing the notations; the master's Software Engineering course was the first proper teaching. By then a homemade process was already in place: underline requirements, pencil sketch, pen. Teacher examples mainly taught notation, not the expected answer.

\textit{Why it matters.} Starting without a taught path (CAT1) can happen as a capstone requirement, not only as a first lab freeze. Later examples still help, but they arrive after students have already invented a sketching routine (CAT4).
\end{tcolorbox}
\vspace{2pt}
\begin{tcolorbox}[colback=gray!4,colframe=black!35,boxrule=0.4pt,left=3pt,right=3pt,top=2pt,bottom=2pt,arc=1pt,title={\small\textbf{L1M1} --- breaking the ice is the main teaching hurdle}]
\footnotesize \textbf{Case:} L1. \textbf{Codes:} L1-C03, L1-C04, L1-C07, L1-C11. \textbf{Turns:} T012, T014, T022, T024, T037.

\textit{Idea.} L1 ranks "breaking the ice"---actually starting to apply modelling to solve a problem---as harder than absorbing semantics alone. He sees most students wait for the teacher to start the solution; only a minority begin independently. Early work involves many typical notation errors until students practise mapping requirements into a model.

\textit{Why it matters.} Lecturer interpretation centres procedural start-up, not just concept difficulty. This is a pedagogy--culture--confidence bundle: students, teaching moves, and expected teacher role interlock before tool or assessment details dominate.
\end{tcolorbox}
\vspace{2pt}
\begin{tcolorbox}[colback=gray!4,colframe=black!35,boxrule=0.4pt,left=3pt,right=3pt,top=2pt,bottom=2pt,arc=1pt,title={\small\textbf{S1M1} --- notation and tool overload block starting}]
\footnotesize \textbf{Case:} S1. \textbf{Codes:} S1-C07, S1-C08, S1-C12. \textbf{Turns:} T018, T020, T029, T051.

\textit{Idea.} On the FreeCol project, dense notation and a feature-rich tool stacked together. Many symbol choices created doubt about the "proper" arrow or relation before any modelling could begin. StarUML's interface felt like too much at once for simple entity-to-box work.

\textit{Why it matters.} STGT: the difficulty is not cognition alone but a socio-technical bundle---course task, UML notation, and StarUML options jointly raised the entry cost. Tool design and notation load shape whether students can start at all.
\end{tcolorbox}
\vspace{2pt}
\begin{tcolorbox}[colback=gray!4,colframe=black!35,boxrule=0.4pt,left=3pt,right=3pt,top=2pt,bottom=2pt,arc=1pt,title={\small\textbf{S3M3} --- notation detail stalls otherwise smooth progress}]
\footnotesize \textbf{Case:} S3. \textbf{Codes:} S3-C10, S3-C11, S3-C12. \textbf{Turns:} T025, T027.

\textit{Idea.} S3 describes BPMN as broadly intuitive until notation details appear: duplicate gateway symbols, expanded subprocesses, unfamiliar constructs, or visually dense diagrams where sequence flow is hard to track. Progress can feel fine until a new symbol or layout problem stops work and triggers re-questioning.

\textit{Why it matters.} Difficulty is not only ``learning BPMN'' once---it recurs at notation boundaries and at the visual/tool layer when models grow. This aligns with a socio-technical view where symbol set, diagram layout, and reading complexity interact.
\end{tcolorbox}
\vspace{2pt}
\begin{tcolorbox}[colback=gray!4,colframe=black!35,boxrule=0.4pt,left=3pt,right=3pt,top=2pt,bottom=2pt,arc=1pt,title={\small\textbf{S2M3} --- framework burden in course not elsewhere}]
\footnotesize \textbf{Case:} S2. \textbf{Codes:} S2-C08, S2-C09, S2-C30, S2-C07. \textbf{Turns:} T022, T024, T026, T103, T106.

\textit{Idea.} Course modelling felt frustrating because unfamiliar frameworks had to be learned before drawing, and abstractions had to fit FreeCol's existing rules. At work and in the dissertation, informal intuitive diagrams sufficed; prescribed frameworks and their multiplicity are off-putting.

\textit{Why it matters.} STGT: the same person experiences modelling differently across contexts---course (framework + legacy codebase) versus self-directed work (purpose-built sketches). Institutional notation requirements may deter practice even when the student values abstraction.
\end{tcolorbox}
\vspace{2pt}
\begin{tcolorbox}[colback=gray!4,colframe=black!35,boxrule=0.4pt,left=3pt,right=3pt,top=2pt,bottom=2pt,arc=1pt,title={\small\textbf{S6M2} --- justification replaces the compile button}]
\footnotesize \textbf{Case:} S6. \textbf{Codes:} S6-C05, S6-C09, S6-C15. \textbf{Turns:} T013, T021, T037.

\textit{Idea.} S6 experiences requirements modelling as prolonged ambiguity: there is no compile button, so correctness comes from walking through the model and explaining why each element exists. The course required justification for actors, dependencies, and goal-to-class mappings; vague briefs and role-play interviews showed requirements shift with who you ask. Initial stress from indefiniteness later read as intentional pedagogy.

\textit{Why it matters.} STGT highlights how evaluation norms (justification, traceability, stakeholder interviews) constitute the "technical" environment alongside tools. Modelling discomfort is socio-epistemic---living without executable feedback---not merely unfamiliar syntax.
\end{tcolorbox}
\vspace{2pt}
\begin{tcolorbox}[colback=gray!4,colframe=black!35,boxrule=0.4pt,left=3pt,right=3pt,top=2pt,bottom=2pt,arc=1pt,title={\small\textbf{S4M2} --- worked examples unlock reasoning and granularity}]
\footnotesize \textbf{Case:} S4. \textbf{Codes:} S4-C03, S4-C04, S4-C05, S4-C06. \textbf{Turns:} T009, T011, T013.

\textit{Idea.} The student can start from text alone but feels less confident without a teacher example, especially when task style is ambiguous. Examples that explain relation choices and contrast too-detailed versus too-abstract diagrams teach process, not just final shapes. A personal quality checklist grew from draft feedback rather than an official rubric.

\textit{Why it matters.} Pedagogical artifacts (explained examples, feedback loops) scaffold judgment about "good enough"---a teaching-tool-learner interplay central to STGT.
\end{tcolorbox}
\vspace{2pt}
\begin{tcolorbox}[colback=gray!4,colframe=black!35,boxrule=0.4pt,left=3pt,right=3pt,top=2pt,bottom=2pt,arc=1pt,title={\small\textbf{S10M5} --- exclusive notation ownership split the system}]
\footnotesize \textbf{Case:} S10. \textbf{Codes:} S10-C01, S10-C12, S10-C13. \textbf{Turns:} T005, T047, T049, T051, T057.

\textit{Idea.} Splitting BPMN / requirements / UML by person produced two versions of the same system. The group invented draft-plus-review and a short carry-over talk before the next model, because the defence asked people individually. Explainability became a practical quality test: if the three authors could not walk it, it was too loaded. Thesis work later dropped strict UML in favour of communicating the architecture.

\textit{Why it matters.} Same split-by-notation risk S7 feared; S10 actually saw divergence, then built the peer-review fix S7 described. CAT9/CAT12: oral defence plus a real reader (thesis) pull models toward simpler, explainable diagrams rather than notation completeness.
\end{tcolorbox}
\vspace{2pt}
\begin{tcolorbox}[colback=gray!4,colframe=black!35,boxrule=0.4pt,left=3pt,right=3pt,top=2pt,bottom=2pt,arc=1pt,title={\small\textbf{S6M1} --- multi-notation RE breaks at translation gaps}]
\footnotesize \textbf{Case:} S6. \textbf{Codes:} S6-C01, S6-C02, S6-C08, S6-C14. \textbf{Turns:} T007, T019, T035.

\textit{Idea.} S6's hardest moment was translating an iStar rationale model into UML: the class diagram looked syntactically fine but lost goal traceability because nobody explained how much strategic content should survive. Each notation family (iStar, NFR, UML) demands a different mindset---organisational/political versus structural---and restarting is expected in a multi-family course, but the missing translation methodology is the real pain. S6 wants one continuous worked example from goals through NFR decisions to UML with documented trace links.

\textit{Why it matters.} Socio-technical RE pedagogy is not only teaching separate notations but orchestrating transitions between them. Without explicit bridge procedures and integrated artefacts, students produce locally correct diagrams that fail globally as requirements models.
\end{tcolorbox}
\vspace{2pt}
\begin{tcolorbox}[colback=gray!4,colframe=black!35,boxrule=0.4pt,left=3pt,right=3pt,top=2pt,bottom=2pt,arc=1pt,title={\small\textbf{S10M1} --- BPMN-to-UML breaks when treated as the same model}]
\footnotesize \textbf{Case:} S10. \textbf{Codes:} S10-C02, S10-C03, S10-C04, S10-C05, S10-C06. \textbf{Turns:} T007, T009, T013, T015, T017, T019, T021.

\textit{Idea.} S10's stuck moment was not a missing symbol. A payments service was clear as an external BPMN pool; a too-direct translation into UML left no place for it (actor, class, component, or omit). The teacher's after-class point was that models answer different questions. Names and requirements carried; gateways, lanes, and tasks did not become classes. One-to-one mapping produced action-classes; reasoning had to restart even though domain knowledge stayed. Traceability was mentioned, not practised across models.

\textit{Why it matters.} This densifies CAT14 at IST AMS beyond S7's syntax-friction account, in the same family as S6's lost iStar-to-UML trace: locally sensible diagrams fail at the bridge. S10 names the missing method in student language (``different questions''), not only as a later teaching wish.
\end{tcolorbox}
\vspace{2pt}
\begin{tcolorbox}[colback=gray!4,colframe=black!35,boxrule=0.4pt,left=3pt,right=3pt,top=2pt,bottom=2pt,arc=1pt,title={\small\textbf{S7M5} --- new syntax at the switch; a real reader later}]
\footnotesize \textbf{Case:} S7. \textbf{Codes:} S7-C01, S7-C14, S7-C15. \textbf{Turns:} T005, T033, T039, T041.

\textit{Idea.} S7 did not call modelling difficult. Confusion was duplicated teaching of the same idea (Databases vs class diagrams) and syntax when moving families (relation types). Once already modelling, the next family was easier to start. Coursework models felt like they only passed through a teacher who had seen millions; the thesis architecture for supervisor and defence became a priority. After graduation S7 modelled to explain ideas without worrying about notation rules, and would rather teaching stressed usefulness and fewer families.

\textit{Why it matters.} CAT14 in this case is light syntax friction plus duplicated course notations, not S6's lost goal-to-UML trace. CAT10 is a clean contrast: same student, low-stakes lecturer audience versus a reader who needed the system explained. Informal workplace modelling is S7's own report, not a lecturer claim.
\end{tcolorbox}
\vspace{2pt}
\begin{tcolorbox}[colback=gray!4,colframe=black!35,boxrule=0.4pt,left=3pt,right=3pt,top=2pt,bottom=2pt,arc=1pt,title={\small\textbf{S1M2} --- worked examples bridge text and diagram}]
\footnotesize \textbf{Case:} S1. \textbf{Codes:} S1-C13, S1-C23, S1-C15. \textbf{Turns:} T031, T033, T071, T073.

\textit{Idea.} S1 learns text-to-diagram mapping by trial and error, starting from components then placing actors. Worked examples---"this description translates to this model"---are named as the most useful feedback form. Templates for common steps would also help.

\textit{Why it matters.} Without explicit process teaching (see S1-C15), example-driven learning becomes the main scaffold. STGT: pedagogy (examples, templates) substitutes for missing procedural instruction and couples learner strategy to course feedback design.
\end{tcolorbox}
\vspace{2pt}
\begin{tcolorbox}[colback=gray!4,colframe=black!35,boxrule=0.4pt,left=3pt,right=3pt,top=2pt,bottom=2pt,arc=1pt,title={\small\textbf{S7M1} --- board labs split starters from copiers}]
\footnotesize \textbf{Case:} S7. \textbf{Codes:} S7-C02, S7-C03, S7-C04, S7-C05. \textbf{Turns:} T009, T011, T013, T017.

\textit{Idea.} IST practicals unlocked work with a simple board example, a class-built requirements list, then a student at the board while the teacher stayed quiet until a student debate finished. S7 started immediately to have something if randomly picked, and because not trying drew more teacher displeasure than a wrong attempt. Classmates often waited and copied the board.

\textit{Why it matters.} The same lab ritual produces two student strategies: attack-to-survive versus wait-to-copy. Lecturer ``students wait to start'' claims can be true of the room without being true of every student.
\end{tcolorbox}
\vspace{2pt}
\begin{tcolorbox}[colback=gray!4,colframe=black!35,boxrule=0.4pt,left=3pt,right=3pt,top=2pt,bottom=2pt,arc=1pt,title={\small\textbf{L4M6} --- active discovery pedagogy resists passivity}]
\footnotesize \textbf{Case:} L4. \textbf{Codes:} L4-C21, L4-C22, L4-C25, L4-C32. \textbf{Turns:} T046, T048, T055, T067.

\textit{Idea.} The lecturer refuses to front-load solved examples, uses icebreakers so shy students cannot wait all semester, and builds solutions gradually from photographed board work while admitting not knowing everything. That contrasts with students overwhelmed by continuous evaluation who rush task-to-task without assimilating.

\textit{Why it matters.} Teaching stance (discovery, openness, pacing) interacts with course load and classroom social format---relevant when shy students prefer private feedback elsewhere.
\end{tcolorbox}
\vspace{2pt}
\begin{tcolorbox}[colback=gray!4,colframe=black!35,boxrule=0.4pt,left=3pt,right=3pt,top=2pt,bottom=2pt,arc=1pt,title={\small\textbf{S5M3} --- early sharing and self-initiated drafts replace waiting}]
\footnotesize \textbf{Case:} S5. \textbf{Codes:} S5-C04, S5-C07, S5-C14. \textbf{Turns:} T009, T017, T033.

\textit{Idea.} S5 rarely waits for a teacher demo: they draft version one, share work-in-progress in lab and group chat, and treat peer feedback as the safety net. Early public sharing reduces anxiety because errors surface early---though S5 acknowledges this pressure could harm shy peers. They want more assessed weekly diagram milestones with peer comments to institutionalise what their team already does informally.

\textit{Why it matters.} Socio-technical conditions (assessment timing, lab culture, chat channels) shape modelling pedagogy as much as notation teaching. A student who initiates drafts and shares early experiences a different course than one who waits for authoritative solutions.
\end{tcolorbox}
\vspace{2pt}
\begin{tcolorbox}[colback=gray!4,colframe=black!35,boxrule=0.4pt,left=3pt,right=3pt,top=2pt,bottom=2pt,arc=1pt,title={\small\textbf{S1M5} --- end-only feedback lets errors accumulate}]
\footnotesize \textbf{Case:} S1. \textbf{Codes:} S1-C22, S1-C18. \textbf{Turns:} T065, T067, T069, T043.

\textit{Idea.} Modelling feedback came mainly at the final project discussion, not iteratively. Problems piled up until they surfaced as "completely wrong" at the end. S1 cannot recall formal validation rules from the course.

\textit{Why it matters.} STGT: feedback timing and assessment design shape learning loops. Without mid-course modelling checks or memorable quality criteria, students rely on self-judgment and peer/dissertation comparisons instead of course scaffolding.
\end{tcolorbox}
\vspace{2pt}
\begin{tcolorbox}[colback=gray!4,colframe=black!35,boxrule=0.4pt,left=3pt,right=3pt,top=2pt,bottom=2pt,arc=1pt,title={\small\textbf{S7M3} --- graded then told: rushed teammate forced a full redo}]
\footnotesize \textbf{Case:} S7. \textbf{Codes:} S7-C09, S7-C10. \textbf{Turns:} T023, T025.

\textit{Idea.} The two-phase AMS project gave feedback after a grade, with no resubmit of that delivery---only a chance to fix the next phase or the defence. S7's group then had to rebuild the first-phase system from scratch. A ``get-it-done-fast'' teammate did almost everything immediately; the two who wanted to plan inherited a rushed model and lost the usual benefit of splitting load.

\textit{Why it matters.} Feedback timing and group tempo interact. Staged deliveries do not create a learning loop if the grade is already locked and a fast member can freeze a weak model before others work.
\end{tcolorbox}
\vspace{2pt}
\begin{tcolorbox}[colback=gray!4,colframe=black!35,boxrule=0.4pt,left=3pt,right=3pt,top=2pt,bottom=2pt,arc=1pt,title={\small\textbf{S10M3} --- timed local rewrite taught a boundary self-check}]
\footnotesize \textbf{Case:} S10. \textbf{Codes:} S10-C09, S10-C10. \textbf{Turns:} T033, T035, T037, T039.

\textit{Idea.} Phase feedback caught a mixed BPMN perspective: an extra pool for something inside the organisation, and message flows that should have been sequence flows. With time left before the next delivery, the group did a large local rewrite, not a full abandon. After that, S10 started asking who is external before drawing. Notation nags were memorised; conceptual comments changed later work.

\textit{Why it matters.} Unlike S7's grade-then-feedback full redo, S10 had time to act and the error was structural, not a rushed teammate. CAT5 is not only ``too late'': when a loop exists, conceptual feedback can become a new start habit.
\end{tcolorbox}
\vspace{2pt}
\begin{tcolorbox}[colback=gray!4,colframe=black!35,boxrule=0.4pt,left=3pt,right=3pt,top=2pt,bottom=2pt,arc=1pt,title={\small\textbf{S4M3} --- shy students need private feedback channels}]
\footnotesize \textbf{Case:} S4. \textbf{Codes:} S4-C07, S4-C08. \textbf{Turns:} T015, T017, T033.

\textit{Idea.} Design ambiguity (multiplicity, class vs attribute) is resolved in office hours or email, rarely in front of the class. Public presentation of half-finished diagrams feels stressful; quick private feedback would enable earlier sharing. The student explicitly asks for complete worked examples in materials so preparation can happen privately before targeted questions.

\textit{Why it matters.} Social format of courses---public vs private validation---shapes modelling participation as much as notation instruction.
\end{tcolorbox}
\vspace{2pt}
\begin{tcolorbox}[colback=gray!4,colframe=black!35,boxrule=0.4pt,left=3pt,right=3pt,top=2pt,bottom=2pt,arc=1pt,title={\small\textbf{S8M2} --- ungraded practice makes checking optional}]
\footnotesize \textbf{Case:} S8. \textbf{Codes:} S8-C02, S8-C05, S8-C06, S8-C07. \textbf{Turns:} T007, T009, T015, T019, T021.

\textit{Idea.} Board examples plus a practice list were not submitted or checked; leaving was allowed. S8 says classmates skipped teacher checks from lack of interest because the exercises did not count, not from shyness. Camunda and Archi were available and unused; draw.io was enough. In the individual project the teacher walked around and almost always asked for changes, but S8 remembers no full redo. Sharing was easy with close friends and with the teacher, not with course-group strangers.

\textit{Why it matters.} Feedback access (CAT6) here is graded-work gated: initiative plus whether the task counts. Optional semantic tools do not get used when a drawing tool already satisfies ``requirements covered.''
\end{tcolorbox}
\vspace{2pt}
\begin{tcolorbox}[colback=gray!4,colframe=black!35,boxrule=0.4pt,left=3pt,right=3pt,top=2pt,bottom=2pt,arc=1pt,title={\small\textbf{S10M2} --- reserved starter: notebook yes, board no}]
\footnotesize \textbf{Case:} S10. \textbf{Codes:} S10-C07, S10-C08, S10-C14. \textbf{Turns:} T025, T027, T029, T031, T055.

\textit{Idea.} S10 solved the exercise in the notebook, then compared with the board. Silence was not ``I cannot start'' and not mainly fear of being wrong; it was dislike of twenty people watching while still thinking. Small arrow-type doubts often stayed unresolved until a classmate, a later search, or the correction. Desk walk-arounds made showing work easy; raising a hand in front of the class did not.

\textit{Why it matters.} Answers S7M1: this reserved IST student is a starter, not a copier. CAT6 here is public-format gating. The same lab that S7 used for teacher goodwill is, for S10, a place to stay quiet unless the teacher comes to the desk.
\end{tcolorbox}
\vspace{2pt}
\begin{tcolorbox}[colback=gray!4,colframe=black!35,boxrule=0.4pt,left=3pt,right=3pt,top=2pt,bottom=2pt,arc=1pt,title={\small\textbf{S4M5} --- tool semantics and workplace informality diverge from class}]
\footnotesize \textbf{Case:} S4. \textbf{Codes:} S4-C01, S4-C12, S4-C13. \textbf{Turns:} T005, T023, T025, T031.

\textit{Idea.} ChatGPT helps explain diagrams in plain language but is rejected for generating graded models. Semantic lab tools catch invalid UML; fast sketch tools allow mistakes at home. A short internship used whiteboard sketches and prototypes---not formal UML---surprising expectations built from coursework.

\textit{Why it matters.} Tool affordances and workplace norms sit alongside classroom notation teaching; misalignment may reinforce "modelling feels pointless" when industry practice looks informal.
\end{tcolorbox}
\vspace{2pt}
\begin{tcolorbox}[colback=gray!4,colframe=black!35,boxrule=0.4pt,left=3pt,right=3pt,top=2pt,bottom=2pt,arc=1pt,title={\small\textbf{S5M4} --- AI and live tools in the group modelling loop}]
\footnotesize \textbf{Case:} S5. \textbf{Codes:} S5-C10, S5-C11, S5-C12. \textbf{Turns:} T023, T025, T027.

\textit{Idea.} S5 uses ChatGPT to brainstorm alternative diagrams and simplify overcomplicated sequence descriptions, then rewrites everything manually before submission. Miro's live co-editing lowered the barrier for group modelling; Visual Paradigm added stricter semantics before coding. They want semantic consistency checking ("this message does not match any operation") more than generative box-drawing. Team passivity, not domain excitement, determines whether modelling effort feels worthwhile.

\textit{Why it matters.} Tools, AI assistants, and team social dynamics form a coupled loop: the same Miro board that enables coordination also needs manual sync when code outpaces diagrams. AI enters as a private brainstorming layer, not a submission shortcut---pointing toward checker-style support rather than replacement modelling.
\end{tcolorbox}
\vspace{2pt}
\begin{tcolorbox}[colback=gray!4,colframe=black!35,boxrule=0.4pt,left=3pt,right=3pt,top=2pt,bottom=2pt,arc=1pt,title={\small\textbf{S10M4} --- AI as a distrusted classmate, not a homework generator}]
\footnotesize \textbf{Case:} S10. \textbf{Codes:} S10-C11, S10-C18. \textbf{Turns:} T041, T043, T045, T057, T059.

\textit{Idea.} S10 used AI after a self-check: screenshot plus system context, small questions (external vs actor), what should carry from BPMN, which model next. Never ``make the whole model.'' It unblocked small doubts and cut some teacher questions, at the cost of missing discussions that taught more.

\textit{Why it matters.} First Portuguese student incident for CAT13, matching S4--S6 explain-not-generate. Here AI also stands in for the public-ask S10 avoids (CAT6).
\end{tcolorbox}
\vspace{2pt}
\begin{tcolorbox}[colback=gray!4,colframe=black!35,boxrule=0.4pt,left=3pt,right=3pt,top=2pt,bottom=2pt,arc=1pt,title={\small\textbf{S3M4} --- Bizagi scaffolds vocabulary and syntax}]
\footnotesize \textbf{Case:} S3. \textbf{Codes:} S3-C14, S3-C15, S3-C16. \textbf{Turns:} T034, T036, T038, T040.

\textit{Idea.} Bizagi Modeler supported learning through in-interface vocabulary, hover explanations, and workflow validity checks. S3 found it intuitive from first open and says it clarified flow distinctions. At scale, one-by-one construction and occasional save corruption made the same tool frustrating for thesis-sized diagrams.

\textit{Why it matters.} Tool affordances are part of the teaching-learning system: they lower vocabulary and syntax entry costs for beginners but do not equally support productivity or reliability at larger model sizes. Tool choice is tied to course reuse across three BPMN contexts and thesis work.
\end{tcolorbox}
\vspace{2pt}
\begin{tcolorbox}[colback=gray!4,colframe=black!35,boxrule=0.4pt,left=3pt,right=3pt,top=2pt,bottom=2pt,arc=1pt,title={\small\textbf{S7M2} --- labs are gut feeling; exams are going on faith}]
\footnotesize \textbf{Case:} S7. \textbf{Codes:} S7-C06, S7-C07, S7-C08. \textbf{Turns:} T019, T021, T035.

\textit{Idea.} Students treated the teacher's pre-drawn board solution as the most correct answer even while the teacher called alternatives a fun fact. In practicals, ``good enough'' was gut feeling and a tick-off of a homemade requirements list. In tests there was no compile signal, so completeness was fear plus faith. An open-source checker was shown once and dropped because the interface created more problems than it solved; remembered tools were family-specific, weak UI, and mainly for pretty PDF export.

\textit{Why it matters.} Assessment risk, not notation knowledge, changes how students judge models. Wanting a compile button (CAT7) can coexist with rejecting the actual checker when tool friction (CAT8) outweighs the check.
\end{tcolorbox}
\vspace{2pt}
\begin{tcolorbox}[colback=gray!4,colframe=black!35,boxrule=0.4pt,left=3pt,right=3pt,top=2pt,bottom=2pt,arc=1pt,title={\small\textbf{L4M5} --- whiteboard compensates for missing integrated tools}]
\footnotesize \textbf{Case:} L4. \textbf{Codes:} L4-C16, L4-C26, L4-C27, L4-C28, L4-C29. \textbf{Turns:} T039, T057, T061, T062, T064.

\textit{Idea.} Without affordable, integrated, semantics-aware tooling across goal-oriented and object-oriented families, the lecturer relies heavily on whiteboard discussion, manual consistency checks, and family-specific tools like piStar for iStar only. Switching tools mid-course creates friction; the top realistic improvement would be a lightweight validator that bridges paradigms.

\textit{Why it matters.} Tool ecology (licenses, semantics, integration) directly constrains what can be taught, checked, and iterated in class---a core STGT pairing of technique and infrastructure.
\end{tcolorbox}
\vspace{2pt}
\begin{tcolorbox}[colback=gray!4,colframe=black!35,boxrule=0.4pt,left=3pt,right=3pt,top=2pt,bottom=2pt,arc=1pt,title={\small\textbf{S1M4} --- simpler tools reduce option overwhelm}]
\footnotesize \textbf{Case:} S1. \textbf{Codes:} S1-C20, S1-C08. \textbf{Turns:} T049, T051, T053, T055.

\textit{Idea.} StarUML made modelling harder through option overload; draw.io felt easier because fewer choices were available. Tool design is seen as capable of improving the modelling experience, separate from notation learning.

\textit{Why it matters.} STGT links tool affordances to learner engagement: a professional-feeling tool can increase cognitive load for students who are not deeply invested in software engineering. Tool simplification is a socio-technical lever alongside teaching.
\end{tcolorbox}
\vspace{2pt}
\begin{tcolorbox}[colback=gray!4,colframe=black!35,boxrule=0.4pt,left=3pt,right=3pt,top=2pt,bottom=2pt,arc=1pt,title={\small\textbf{L1M5} --- proper tools versus draw.io in teaching stance}]
\footnotesize \textbf{Case:} L1. \textbf{Codes:} L1-C17, L1-C18, L1-C20, L1-C21. \textbf{Turns:} T057, T059, T061, T063, T071, T073.

\textit{Idea.} L1 steers students away from draw.io toward metamodel-backed tools (e.g. Visual Paradigm) but struggles with licensing, imperfect checking, and past bad tools he abandoned. Class time is too short for tool onboarding and troubleshooting. He enthusiastically endorses a hypothetical faculty-provided free tool with feedback, checklists, and next-step guidance---explicitly a design preference, not an implemented system.

\textit{Why it matters.} Tool policy is part of teaching: what counts as "modelling" vs drawing, what errors are caught, and who pays (institution vs student). Lecturer ideal tool features partially overlap student reports (vocabulary help, syntax warnings) but add institutional provisioning.
\end{tcolorbox}
\vspace{2pt}
\begin{tcolorbox}[colback=gray!4,colframe=black!35,boxrule=0.4pt,left=3pt,right=3pt,top=2pt,bottom=2pt,arc=1pt,title={\small\textbf{S4M1} --- group work can kill modelling when code wins}]
\footnotesize \textbf{Case:} S4. \textbf{Codes:} S4-C02, S4-C09, S4-C11. \textbf{Turns:} T007, T017, T021.

\textit{Idea.} In a generic shop group project, teammates pushed Spring implementation while the student maintained the class diagram alone; dismissal as "paperwork" stopped refinement until the model no longer matched code. Modelling worked later when one person owned an interesting library design; early sharing in groups turned diagrams into arguments rather than improvements.

\textit{Why it matters.} Socio-technical conditions---ownership split, domain interest, group dynamics---can matter more than UML notation knowledge for whether modelling persists.
\end{tcolorbox}
\vspace{2pt}
\begin{tcolorbox}[colback=gray!4,colframe=black!35,boxrule=0.4pt,left=3pt,right=3pt,top=2pt,bottom=2pt,arc=1pt,title={\small\textbf{L4M3} --- domain freedom trades lecturer workload for fairness}]
\footnotesize \textbf{Case:} L4. \textbf{Codes:} L4-C06, L4-C07, L4-C08, L4-C11. \textbf{Turns:} T026, T027, T028, T033.

\textit{Idea.} Letting groups choose or refine domains---scoped when too simple or complex---increases marking and feedback work but avoids imposing uninterested domains the lecturer remembers as unfair. Modelling is reserved for problems complex enough to need communication beyond a weekend build. Under semester overload, some students still ask for tighter "just tell me what to do" tasks.

\textit{Why it matters.} Course conditions (freedom, fairness, workload, concurrent evaluations) co-shape whether modelling feels worth the effort for both teacher and learners.
\end{tcolorbox}
\vspace{2pt}
\begin{tcolorbox}[colback=gray!4,colframe=black!35,boxrule=0.4pt,left=3pt,right=3pt,top=2pt,bottom=2pt,arc=1pt,title={\small\textbf{L1M3} --- missing theory class breaks slide-only catch-up}]
\footnotesize \textbf{Case:} L1. \textbf{Codes:} L1-C14, L1-C15, L1-C22. \textbf{Turns:} T047, T051, T077.

\textit{Idea.} L1 reports students skipping theoretical sessions, then learning from slides alone, misinterpreting concepts, and arriving at practicals unprepared. He does not repeat theory in labs and resists making slides a full tutorial. He attributes absences partly to Portuguese student culture, heavy multi-course workload, and continuous evaluation pressure---not lack of interest in modelling alone.

\textit{Why it matters.} Course structure assumes theory attendance; practicals "just use" concepts. Socio-organisational conditions (semester load, evaluation design) undermine that assumption and create a theory--practice gap independent of notation choice.
\end{tcolorbox}
\vspace{2pt}
\begin{tcolorbox}[colback=gray!4,colframe=black!35,boxrule=0.4pt,left=3pt,right=3pt,top=2pt,bottom=2pt,arc=1pt,title={\small\textbf{S7M6} --- competing courses made modelling a task to despachar}]
\footnotesize \textbf{Case:} S7. \textbf{Codes:} S7-C16. \textbf{Turns:} T013, T035.

\textit{Idea.} S7 called the modelling workload adequate, but a packed second-year semester plus low interest meant finishing those tasks fast to protect other courses. Drawing and modelling were not a strength or a focus. That condition sits beside the board-lab split: some students only copied, and S7 still started---partly for the teacher's good books.

\textit{Why it matters.} CAT11 here is interest plus concurrent load, not an overloaded modelling course itself. Student ``despachar'' is not the same as lecturers' ECTS-scarcity story.
\end{tcolorbox}
\vspace{2pt}
\begin{tcolorbox}[colback=gray!4,colframe=black!35,boxrule=0.4pt,left=3pt,right=3pt,top=2pt,bottom=2pt,arc=1pt,title={\small\textbf{S8M3} --- similar families confuse; paid work drains care}]
\footnotesize \textbf{Case:} S8. \textbf{Codes:} S8-C08, S8-C09. \textbf{Turns:} T023, T025, T027, T031.

\textit{Idea.} Switching was hardest when diagrams looked alike (requirements vs class relations); an example beside the task made it fine, and knowing one family made the next easier. Models still felt like drawings next to more tangible code, so perfection mattered less than presence. A tutoring job in the master's, not competing courses, cut interest. S8 now doubts professional utility and would teach workplace benefits rather than more notation.

\textit{Why it matters.} CAT14 here is lookalike-notation confusion, not S6's lost traceability. CAT11's competing load is a paid job, not ECTS stacking. The utility wish is a student interpretation, not evidence that modelling ``will die.''
\end{tcolorbox}
\vspace{2pt}
